\documentclass[11pt]{article}

\usepackage[margin=1in]{geometry}
\usepackage{amsmath, amssymb, amsthm}
\usepackage{parskip}
\usepackage{graphicx}
\usepackage{float}

\begin{document}

\begin{center}
{\Large \textbf{MECH 309: Assignment 6, Question 2}}\\
Cagri Arslan\\
\today
\end{center}

Consider the initial value problem
\begin{equation}
\dot{x}(t)=f(x(t)), \qquad x(0)=x_0.
\end{equation}
We derive numerical integration schemes using finite-difference approximations to the derivative.

\section*{(a) Forward difference and Forward Euler}

Recall the forward-difference approximation
\begin{equation}
\dot{x}(t_k) \approx \frac{x_{k+1}-x_k}{h},
\end{equation}
where $h=t_{k+1}-t_k$.

Evaluating the right-hand side of the ODE at $x_k=x(t_k)$ gives
\begin{equation}
\dot{x}(t_k)=f(x_k).
\end{equation}
Therefore,
\begin{equation}
\frac{x_{k+1}-x_k}{h}=f(x_k).
\end{equation}
Rearranging,
\begin{equation}
\boxed{x_{k+1}=x_k+h\,f(x_k).}
\end{equation}
This is the \textbf{Forward Euler} method.

To determine the accuracy, expand $x(t_k+h)$ about $t_k$:
\begin{equation}
x(t_k+h)=x(t_k)+h\dot{x}(t_k)+\frac{h^2}{2}\ddot{x}(t_k)+\frac{h^3}{6}x^{(3)}(t_k)+\cdots.
\end{equation}
Subtracting $x(t_k)$ and dividing by $h$ yields
\begin{equation}
\frac{x(t_k+h)-x(t_k)}{h}=\dot{x}(t_k)+\frac{h}{2}\ddot{x}(t_k)+O(h^2).
\end{equation}
Hence,
\begin{equation}
\dot{x}(t_k)=\frac{x_{k+1}-x_k}{h}+O(h),
\end{equation}
so the Forward Euler method is \textbf{first-order accurate}, i.e. \(O(h)\).

\section*{(b) Backward difference and Backward Euler}

Recall the backward-difference approximation
\begin{equation}
\dot{x}(t_k) \approx \frac{x_k-x_{k-1}}{h},
\end{equation}
where $h=t_k-t_{k-1}$.

Evaluating the ODE at $x_k=x(t_k)$ gives
\begin{equation}
\dot{x}(t_k)=f(x_k).
\end{equation}
Thus,
\begin{equation}
\frac{x_k-x_{k-1}}{h}=f(x_k).
\end{equation}
Rearranging,
\begin{equation}
\boxed{x_k=x_{k-1}+h\,f(x_k).}
\end{equation}
Writing this one step forward gives the update form
\begin{equation}
\boxed{x_{k+1}=x_k+h\,f(x_{k+1}).}
\end{equation}
This is the \textbf{Backward Euler} method.

Using the Taylor expansion of $x(t_k-h)$ about $t_k$,
\begin{equation}
x(t_k-h)=x(t_k)-h\dot{x}(t_k)+\frac{h^2}{2}\ddot{x}(t_k)-\frac{h^3}{6}x^{(3)}(t_k)+\cdots,
\end{equation}
we obtain
\begin{equation}
\frac{x(t_k)-x(t_k-h)}{h}=\dot{x}(t_k)-\frac{h}{2}\ddot{x}(t_k)+O(h^2).
\end{equation}
Hence,
\begin{equation}
\dot{x}(t_k)=\frac{x_k-x_{k-1}}{h}+O(h),
\end{equation}
so Backward Euler is also \textbf{first-order accurate}, i.e. \(O(h)\).

\section*{(c) Central difference and the Leapfrog method}

Recall the central-difference approximation
\begin{equation}
\dot{x}(t_k) \approx \frac{x_{k+1}-x_{k-1}}{2h},
\end{equation}
where $h=t_{k+1}-t_k=t_k-t_{k-1}$.

Evaluating the ODE at $x_k=x(t_k)$ gives
\begin{equation}
\dot{x}(t_k)=f(x_k).
\end{equation}
Therefore,
\begin{equation}
\frac{x_{k+1}-x_{k-1}}{2h}=f(x_k).
\end{equation}
Rearranging,
\begin{equation}
\boxed{x_{k+1}=x_{k-1}+2h\,f(x_k).}
\end{equation}
This is the \textbf{leapfrog method}.

To determine the accuracy, use the Taylor expansions
\begin{align}
x(t_k+h) &= x(t_k)+h\dot{x}(t_k)+\frac{h^2}{2}\ddot{x}(t_k)+\frac{h^3}{6}x^{(3)}(t_k)+O(h^4),\\
x(t_k-h) &= x(t_k)-h\dot{x}(t_k)+\frac{h^2}{2}\ddot{x}(t_k)-\frac{h^3}{6}x^{(3)}(t_k)+O(h^4).
\end{align}
Subtracting gives
\begin{equation}
x(t_k+h)-x(t_k-h)=2h\dot{x}(t_k)+\frac{h^3}{3}x^{(3)}(t_k)+O(h^5).
\end{equation}
Dividing by $2h$,
\begin{equation}
\frac{x(t_k+h)-x(t_k-h)}{2h}=\dot{x}(t_k)+O(h^2).
\end{equation}
Thus the leapfrog method is \textbf{second-order accurate}, i.e. \(O(h^2)\).

To start the leapfrog method, we need both $x_0$ and $x_1$, but the initial condition only provides $x_0$. We therefore compute $x_1$ using a one-step method such as Forward Euler:
\begin{equation}
\boxed{x_1=x_0+h\,f(x_0).}
\end{equation}
Then the leapfrog update can be applied for $k\geq 1$:
\begin{equation}
\boxed{x_{k+1}=x_{k-1}+2h\,f(x_k).}
\end{equation}
For example,
\begin{equation}
x_2=x_0+2h\,f(x_1), \qquad x_3=x_1+2h\,f(x_2), \quad \text{etc.}
\end{equation}

\end{document}
